\documentclass[hyphens,aspectratio=169,dvipsnames]{beamer}
\usepackage{graphicx}
\usepackage{xcolor}
\usepackage{amssymb}
\usepackage{amsmath}
\usepackage{mathtools}
\usepackage{textcomp}
\usepackage{moresize}
\usepackage{framed}
\usepackage{minted}
\usepackage{relsize}
\usepackage{emoji}
\usepackage{tikz}
\usepackage{comment}
\usetikzlibrary{shapes.geometric, arrows, positioning}
\usetheme{Berlin}
\usecolortheme[RGB={0,95,47}]{structure}
\beamertemplatenavigationsymbolsempty

\makeatletter
\AtBeginEnvironment{minted}{\dontdofcolorbox}
\def\dontdofcolorbox{\renewcommand\fcolorbox[4][]{##4}}
\makeatother

\newcommand{\textpf}[1]{\texttt{\color{black}\fcolorbox{lightgray}{lightgray}{#1}}}

\begin{document}

\begin{frame}[fragile]{Topics}
    \begin{itemize}
        \item Shared Memory
        \item Sockets
    \end{itemize}
\end{frame}

\begin{frame}[fragile]{Shared Memory}
    \begin{itemize}
        \pause \item Shared memory is a Unix IPC mechanism that allows two processes to share a memory mapping.
        \pause \item The easiest way to make a shared memory mapping is with \textpf{mmap}.
        \pause \item When you \textpf{mmap} pages with \textpf{MAP\_SHARED} instead of \textpf{MAP\_PRIVATE}, those pages are shared into child processes.
        \pause \item When child processes write into those mappings, the parent can see the writes, and vice-versa.
        \pause \item Demo!
        \pause \item Now with optimizations enabled?
    \end{itemize}
\end{frame}

\begin{frame}[fragile]{\textpf{mmap}ping Files}
    \begin{itemize}
        \pause \item As we've alluded to in previous classes, \textpf{mmap} can be used to map the contents of a file into your address space.
        \pause \item To do this, don't use the \textpf{MAP\_ANONYMOUS} flag, and change the \textpf{fd} argument to a valid file descriptor.
        \pause \item You can use this to create shared mappings between processes that don't have a parent-child relationship!
    \end{itemize}
\end{frame}

\begin{frame}[fragile]{TCP (Background)}
    \begin{itemize}
        \pause \item TCP (transmission control protocol) is a protocol that allows two computers to establish a pipe over the internet.
        \pause \item That is, one computer can write into the pipe, and the other computer can read the written bytes out of the other end of the pipe (and vice-versa).
        \pause \item When you connect to a website in the browser, you're very likely doing so using TCP connections.
    \end{itemize}
\end{frame}

\begin{frame}[fragile]{\textpf{ncat}}
    \begin{itemize}
        \pause \item Netcat (\textpf{nc(1)}) is a command-line utility for creating TCP connections.
        \pause \item Demo! % First use nc as a client to talk to a remote server,
                           % Then as a server and establish a connection using another nc
    \end{itemize}
\end{frame}

\begin{frame}[fragile]{Sockets}
    \begin{itemize}
        \pause \item Sockets are file descriptors that correspond to network connections.
        \pause \item Sockets are really versatile, and support a whole lot of protocols.
        \pause \item Today, we'll be using sockets to establish TCP connections.
        \pause \item (If you care to learn more about this, take CS60).
    \end{itemize}
\end{frame}

\begin{frame}[fragile]{\textpf{socket(2)}}
    \begin{itemize}
        \pause \item You can create a socket using the \textpf{socket} system call.
        \pause \item \textpf{socket} takes 3 arguments:
            \begin{itemize}
                \pause \item \textpf{domain}: The desired protocol family for the socket to use. e.g., \textpf{AF\_INET} for internet protocols, \textpf{AF\_BLUETOOTH} for bluetooth, etc.
                \pause \item \textpf{type}: The desired communication semantics. Use \textpf{SOCK\_STREAM} for streaming protocols like TCP.
                \pause \item \textpf{protocol}: The desired communication protocol. If left to 0, gets a default value determined by \textpf{type}. Use \textpf{IPPROTO\_TCP} (defined in \textpf{netinet/in.h}) for TCP.
            \end{itemize}
        \pause \item \textpf{socket} returns the file descriptor corresponding to your socket.
    \end{itemize}
\end{frame}

\begin{frame}[fragile]{\textpf{connect(2)}}
    \begin{itemize}
        \pause \item You can connect a socket to a remote peer using the \textpf{connect} system call.
        \pause \item \textpf{connect} takes 3 arguments:
            \begin{itemize}
                \pause \item \textpf{int sockfd}: The socket's file descriptor.
                \pause \item \textpf{struct sockaddr const *addr}: A \textpf{struct} representing the address of the host that you want to connect to.
                \pause \item \textpf{socklen\_t addrlen}: The size of \textpf{addr}.
            \end{itemize}
        \pause \item \textpf{connect} returns an error code on failure, and 0 on success.
    \end{itemize}
\end{frame}


\begin{frame}[fragile]{Interacting with a socket}
    Sockets can be read, written, and closed just like any other file descriptor.
\end{frame}

\begin{frame}[fragile]{Activity: Steal the secret!}
    I have started a TCP server on Plink's port 1234 that writes out a secret string.
    
    Connect to this server using \textpf{nc} and print the string.

    Then, write a C program to do the same thing.
\end{frame}

\end{document}
